
% Page 57

Bien que ce soit d'un autre ordre, peut-être, faut-il faire remonter à ce début du XVII\textsuperscript{ème} siècle, ce géant de verdure que constitue le « Hêtre de Margou », communément appelé en patois : « lou faous dé Roussignol ». La chose n'est pas invraisemblable, car sur la fin du XVII\textsuperscript{ème} siècle le cadastre de Paulin (1693) signale que Jean Sendral de Margou a des biens avoisinant des terres et bois et fontaine dits « de Rossignol ». Les Rossignol ont déjà laissé leur nom à des terrains et à l'arbre majestueux qui domine encore et qui est devenu point géodésique. Ce vétéran qui domine encore le sommet de Margou a bien pu naître au moment où la sage administration de Sully favorisait le reboisement du pays et progresser l'agriculture devenue « mamelle de la France ». De ci, de là, en ce temps-là, poussèrent les ormeaux de Sully et l'un d'entre eux qui se refuse à mourir, a bien pu voir le jour sur la petite place de Notre Dame d'Ourtiquet grâce à ce ministre d'Henri IV. On ne dit pas qui fut son parrain comme c'est le cas semble-t-il pour l'arbre phénomène de Margou. Salut à tous ces vestiges d'un passé tourmenté.
